\documentclass[12pt,a4paper]{article}

\usepackage[T1]{fontenc}
\usepackage[utf8]{inputenc}
\usepackage[brazil]{babel}
\usepackage{lmodern}
\usepackage[a4paper,margin=2.5cm]{geometry}
\usepackage{graphicx}
\usepackage{booktabs}
\usepackage{longtable}
\usepackage{float}
\usepackage{amsmath}
\usepackage{hyperref}
\usepackage{xcolor}
\usepackage{enumitem}

\hypersetup{
    colorlinks=true,
    linkcolor=blue!50!black,
    urlcolor=blue!50!black,
    citecolor=blue!50!black
}

\title{Base T\'ecnica do Projeto\\Classifica\c{c}\~ao de Sinais de ECG com MIT-BIH}
\author{Tiago Leite Pereira}
\date{\today}

\begin{document}

\maketitle

\begin{abstract}
Este documento consolida a base t\'ecnica do projeto de classifica\c{c}\~ao de sinais de eletrocardiograma (ECG) desenvolvido sobre a MIT-BIH Arrhythmia Database. O objetivo do texto n\~ao \'e apenas registrar tarefas executadas, mas explicitar a formula\c{c}\~ao do problema, as decis\~oes metodol\'ogicas j\'a implementadas, o fluxo de dados, as representa\c{c}\~oes extra\'idas do sinal e o protocolo de avalia\c{c}\~ao usado at\'e o momento. O documento foi estruturado como base de manuscrito cient\'ifico, de modo que possa evoluir para introdu\c{c}\~ao, metodologia, resultados e discuss\~ao de um artigo completo.
\end{abstract}

\tableofcontents
\newpage

\section{Prop\'osito do manuscrito}
Este texto serve como base t\'ecnica para um futuro artigo cient\'ifico. Por isso, ele foi organizado para responder de forma objetiva quatro perguntas centrais:

\begin{enumerate}[label=\arabic*.]
    \item qual problema est\'a sendo atacado;
    \item quais dados e representa\c{c}\~oes est\~ao sendo usados;
    \item como o pipeline atualmente implementado transforma sinal bruto em decis\~ao de classe;
    \item quais evid\^encias experimentais j\'a existem e quais lacunas permanecem abertas.
\end{enumerate}

Em vez de funcionar apenas como relat\'orio operacional, o documento busca capturar a l\'ogica cient\'ifica do projeto: hip\'oteses impl\'icitas, escolhas de modelagem, protocolo de valida\c{c}\~ao e riscos metodol\'ogicos.

\section{Problema de pesquisa}
O problema geral consiste em classificar batimentos card\'iacos a partir de sinais de ECG, tomando como base registros anotados da MIT-BIH Arrhythmia Database. O reposit\'orio foi organizado para sustentar duas linhas metodol\'ogicas complementares:

\begin{itemize}
    \item uma linha baseada em din\^amica simb\'olica e descritores n\~ao lineares;
    \item uma linha baseada em modelos temporais orientados por dados, com \^enfase inicial em LSTM e extens\~oes do tipo CNN/LSTM.
\end{itemize}

No estado atual do projeto, a trilha mais madura \'e o baseline cl\'assico por batimento, apoiado por atributos morfol\'ogicos, atributos simb\'olicos e descritores de intervalo RR.

\section{Base de dados}
\subsection{MIT-BIH Arrhythmia Database}
O conjunto de dados de refer\^encia utilizado no projeto \'e a MIT-BIH Arrhythmia Database. No reposit\'orio, os dados brutos est\~ao organizados em \texttt{data/raw/mit\_bih/}. Cada registro cont\'em pelo menos:

\begin{itemize}
    \item o tra\c{c}ado do ECG em um ou mais canais;
    \item metadados do registro;
    \item anota\c{c}\~oes de eventos card\'iacos, incluindo a localiza\c{c}\~ao amostral dos batimentos.
\end{itemize}

O pipeline usa taxa de amostragem de 360 Hz, conforme parametriza\c{c}\~ao atual do projeto.

\subsection{Estrutura dos dados no reposit\'orio}
Os dados s\~ao organizados em tr\^es camadas:

\begin{longtable}{p{0.28\textwidth} p{0.64\textwidth}}
\toprule
\textbf{Camada} & \textbf{Fun\c{c}\~ao} \\
\midrule
\texttt{data/raw/mit\_bih/} & arquivos originais da base MIT-BIH em formato WFDB; \\
\texttt{data/interim/signal\_tables/} & vers\~oes tabulares em CSV dos sinais e das anota\c{c}\~oes; \\
\texttt{data/processed/} & sinais processados, features agregadas e datasets derivados. \\
\bottomrule
\end{longtable}

Essa separa\c{c}\~ao favorece reprodutibilidade, inspe\c{c}\~ao intermedi\'aria e reaproveitamento das etapas do pipeline.

\section{Pergunta metodol\'ogica central}
No n\'ucleo do projeto, a seguinte pergunta guia a implementa\c{c}\~ao atual:

\begin{quote}
Como representar batimentos de ECG de forma suficientemente informativa para distinguir padr\~oes fisiol\'ogicos e arritmias, preservando reprodutibilidade experimental e capacidade de compara\c{c}\~ao entre abordagens cl\'assicas e modelos temporais?
\end{quote}

O pipeline hoje implementado responde a essa pergunta decompondo o problema em quatro etapas:

\begin{enumerate}[label=\arabic*.]
    \item prepara\c{c}\~ao e limpeza do sinal;
    \item localiza\c{c}\~ao de batimentos ou uso das anota\c{c}\~oes como refer\^encia estrutural;
    \item extra\c{c}\~ao de descritores por batimento;
    \item treinamento e avalia\c{c}\~ao de classificadores.
\end{enumerate}

\section{Arquitetura do pipeline}
\subsection{Organiza\c{c}\~ao geral}
O pacote principal do projeto est\'a em \texttt{src/ecg\_classification/} e foi dividido em seis blocos:

\begin{itemize}
    \item \texttt{io}: leitura de sinais, anota\c{c}\~oes e exporta\c{c}\~ao para CSV;
    \item \texttt{preprocessing}: limpeza, normaliza\c{c}\~ao e corre\c{c}\~ao de baseline;
    \item \texttt{segmentation}: detec\c{c}\~ao de picos R e utilidades de segmenta\c{c}\~ao;
    \item \texttt{features}: constru\c{c}\~ao de descritores simb\'olicos, morfol\'ogicos e RR;
    \item \texttt{models}: modelos cl\'assicos e scaffolds para redes neurais;
    \item \texttt{evaluation}: m\'etricas, splits por registro e avalia\c{c}\~ao de detec\c{c}\~ao de picos.
\end{itemize}

\subsection{Pontos de entrada}
Os scripts principais atualmente dispon\'iveis s\~ao:

\begin{verbatim}
python scripts/make_dataset.py
python scripts/preprocess_mitbih.py
python scripts/extract_symbolic_features.py
python scripts/train_classical_ml.py
python scripts/evaluate_rpeak_detection.py
\end{verbatim}

O projeto tamb\'em possui um despachante de execu\c{c}\~ao:

\begin{verbatim}
python scripts/run_experiment.py make_dataset
python scripts/run_experiment.py preprocess
python scripts/run_experiment.py symbolic_features
python scripts/run_experiment.py train_classical_ml
python scripts/run_experiment.py train_lstm
\end{verbatim}

\section{Metodologia implementada}
\subsection{Ingest\~ao e tabulariza\c{c}\~ao}
A etapa de ingest\~ao transforma os registros WFDB em tabelas CSV. Para cada registro, o projeto exporta:

\begin{itemize}
    \item um arquivo \texttt{\_record.csv} com os canais do ECG;
    \item um arquivo \texttt{\_annotation.csv} com as colunas de amostra e s\'imbolo anotado.
\end{itemize}

Essa decis\~ao reduz o acoplamento com o formato nativo do WFDB durante as etapas de explora\c{c}\~ao, depura\c{c}\~ao e engenharia de atributos.

\subsection{Pr\'e-processamento do sinal}
O pr\'e-processamento atual segue uma sequ\^encia simples e reproduz\'ivel:

\begin{enumerate}[label=\arabic*.]
    \item sele\c{c}\~ao do canal principal do registro;
    \item limpeza do ECG com \texttt{nk.ecg\_clean} do NeuroKit2;
    \item remo\c{c}\~ao de baseline wander por decomposi\c{c}\~ao wavelet;
    \item salvamento das vers\~oes \texttt{raw}, \texttt{cleaned}, \texttt{baseline} e \texttt{baseline\_corrected}.
\end{enumerate}

Em termos metodol\'ogicos, essa etapa busca reduzir contamina\c{c}\~ao de baixa frequ\^encia e tornar a representa\c{c}\~ao dos batimentos mais est\'avel para segmenta\c{c}\~ao e extra\c{c}\~ao de features.

\subsection{Detec\c{c}\~ao de picos R}
A detec\c{c}\~ao de picos R atualmente \'e implementada com \texttt{nk.ecg\_peaks} do NeuroKit2. O projeto trata essa etapa como componente metodol\'ogico cr\'itico por dois motivos:

\begin{itemize}
    \item erros na localiza\c{c}\~ao do pico deslocam a janela do batimento e degradam as features extra\'idas;
    \item a qualidade da detec\c{c}\~ao afeta diretamente qualquer pipeline que dependa de janelas centradas no batimento.
\end{itemize}

Por isso, foi adicionada uma rotina espec\'ifica de avalia\c{c}\~ao da detec\c{c}\~ao, comparando picos detectados com as anota\c{c}\~oes de refer\^encia da base.

\subsection{Constru\c{c}\~ao do dataset por batimento}
O dataset supervisionado por batimento \'e gerado a partir das anota\c{c}\~oes da base. Para cada batimento v\'alido:

\begin{enumerate}[label=\arabic*.]
    \item identifica-se a amostra anotada como pico de refer\^encia;
    \item extrai-se uma janela com 100 amostras \`a esquerda e 150 \`a direita;
    \item descartam-se casos que ultrapassem as bordas do sinal;
    \item associa-se ao batimento seu r\'otulo e seus descritores.
\end{enumerate}

Essa janela corresponde a 250 amostras por batimento, o que equivale a aproximadamente 694 ms em 360 Hz.

\subsection{Mapeamento de classes}
Antes do treinamento do baseline cl\'assico, os s\'imbolos originais da MIT-BIH s\~ao filtrados para manter apenas r\'otulos de batimento e depois mapeados para o esquema AAMI:

\begin{itemize}
    \item \texttt{N}: batimentos normais e variantes pr\'oximas;
    \item \texttt{S}: batimentos supraventriculares;
    \item \texttt{V}: batimentos ventriculares;
    \item \texttt{F}: fus\~ao;
    \item \texttt{Q}: batimentos n\~ao classific\'aveis ou associados a marca\c{c}\~oes espec\'ificas.
\end{itemize}

Essa padroniza\c{c}\~ao \'e importante para comparabilidade com literatura e para reduzir a fragmenta\c{c}\~ao excessiva das classes.

\subsection{Descritores extra\'idos}
\subsubsection{Descritores morfol\'ogicos}
Para cada janela de batimento, o projeto extrai atualmente cinco descritores morfol\'ogicos simples:

\begin{itemize}
    \item m\'edia;
    \item desvio-padr\~ao;
    \item valor m\'aximo;
    \item valor m\'inimo;
    \item amplitude pico-a-pico.
\end{itemize}

Esses atributos servem como baseline de baixa complexidade e oferecem uma representa\c{c}\~ao resumida da forma do batimento.

\subsubsection{Descritores de din\^amica simb\'olica}
O componente simb\'olico hoje implementado segue a seguinte l\'ogica:

\begin{enumerate}[label=\arabic*.]
    \item normaliza\c{c}\~ao do sinal por z-score;
    \item parti\c{c}\~ao do intervalo normalizado em quatro bins;
    \item simboliza\c{c}\~ao do sinal com base nesses bins;
    \item forma\c{c}\~ao de palavras simb\'olicas de tamanho 3;
    \item c\'alculo de entropia de Shannon e estat\'isticas derivadas.
\end{enumerate}

Os atributos simb\'olicos atualmente computados s\~ao:

\begin{itemize}
    \item entropia simb\'olica;
    \item n\'umero de palavras distintas;
    \item m\'edia dos s\'imbolos;
    \item desvio-padr\~ao dos s\'imbolos.
\end{itemize}

Do ponto de vista cient\'ifico, essa representa\c{c}\~ao busca capturar padr\~oes de complexidade e organiza\c{c}\~ao do tra\c{c}ado para al\'em de estat\'isticas de amplitude.

\subsubsection{Descritores RR}
Para cada batimento, o projeto calcula:

\begin{itemize}
    \item \texttt{rr\_prev}: intervalo at\'e o batimento anterior;
    \item \texttt{rr\_next}: intervalo at\'e o batimento seguinte.
\end{itemize}

Esses atributos incorporam contexto temporal local, que \'e relevante para discriminar arritmias mesmo quando a morfologia isolada do batimento n\~ao \'e suficiente.

\subsection{Modelo baseline atual}
O baseline mais consolidado do reposit\'orio utiliza um pipeline com:

\begin{itemize}
    \item padroniza\c{c}\~ao via \texttt{StandardScaler};
    \item classificador \texttt{BalancedRandomForestClassifier} quando \texttt{imbalanced-learn} est\'a dispon\'ivel;
    \item alternativa com \texttt{RandomForestClassifier} e pondera\c{c}\~ao de classes quando a depend\^encia n\~ao est\'a presente.
\end{itemize}

Foram definidos 300 estimadores, com \^enfase expl\'icita em lidar com desbalanceamento entre classes.

\section{Protocolo de avalia\c{c}\~ao}
\subsection{Split por registro}
O pipeline adota separa\c{c}\~ao por registro, e n\~ao por batimento aleat\'orio. Essa escolha \'e metodologicamente importante porque reduz vazamento de informa\c{c}\~ao entre treino e teste, j\'a que batimentos do mesmo paciente ou registro tendem a compartilhar padr\~oes muito pr\'oximos.

No c\'odigo atual, a divis\~ao est\'a parametrizada como:

\begin{itemize}
    \item treino: 70\%;
    \item valida\c{c}\~ao: 15\%;
    \item teste: 15\%.
\end{itemize}

Na rotina de treinamento do baseline cl\'assico, o desempenho reportado atualmente usa treino e teste por registro.

\subsection{M\'etricas de classifica\c{c}\~ao}
As m\'etricas atualmente computadas para o classificador s\~ao:

\begin{itemize}
    \item \textit{accuracy};
    \item \textit{precision} macro;
    \item \textit{recall} macro;
    \item \textit{F1} macro;
    \item matriz de confus\~ao.
\end{itemize}

O uso de m\'edias macro \'e especialmente relevante neste contexto, pois evita que classes majorit\'arias dominem completamente a leitura do desempenho.

\subsection{M\'etricas para detec\c{c}\~ao de picos}
Para a detec\c{c}\~ao de picos R, o projeto implementa pareamento entre picos detectados e picos anotados com toler\^ancia fixa. Dado um conjunto de verdade de refer\^encia e um conjunto de picos detectados, s\~ao computados:

\begin{align}
\text{Precision} &= \frac{TP}{TP + FP} \\
\text{Recall} &= \frac{TP}{TP + FN} \\
\text{F1} &= \frac{2 \cdot \text{Precision} \cdot \text{Recall}}{\text{Precision} + \text{Recall}}
\end{align}

Tamb\'em s\~ao calculados erro absoluto m\'edio e erro absoluto mediano, em amostras e em milissegundos, para os pares corretamente associados.

\section{Estado experimental atual}
\subsection{Resultados do baseline cl\'assico}
O arquivo \texttt{reports/tables/classical\_ml\_metrics.json} registra, no estado atual do projeto, os seguintes valores:

\begin{table}[H]
\centering
\begin{tabular}{lc}
\toprule
\textbf{M\'etrica} & \textbf{Valor} \\
\midrule
Accuracy & 0.7543 \\
Precision macro & 0.3807 \\
Recall macro & 0.3845 \\
F1 macro & 0.3713 \\
\bottomrule
\end{tabular}
\caption{Desempenho atual do baseline cl\'assico por batimento.}
\end{table}

Esses n\'umeros indicam que o classificador j\'a consegue capturar parte do padr\~ao global, mas ainda apresenta limita\c{c}\~oes importantes quando observado sob m\'etricas macro. Em outras palavras, o \textit{accuracy} isolado n\~ao descreve adequadamente o desempenho em classes minorit\'arias.

\subsection{Resultados da avalia\c{c}\~ao de picos R}
A nova rotina de avalia\c{c}\~ao da detec\c{c}\~ao de picos produziu os seguintes valores m\'edios globais, com toler\^ancia de 15 amostras:

\begin{table}[H]
\centering
\begin{tabular}{lc}
\toprule
\textbf{M\'etrica} & \textbf{Valor} \\
\midrule
Precision & 0.9490 \\
Recall & 0.9581 \\
F1 & 0.9523 \\
\bottomrule
\end{tabular}
\caption{M\'edias globais da detec\c{c}\~ao de picos R.}
\end{table}

Os registros com pior desempenho observado at\'e o momento foram:

\begin{table}[H]
\centering
\begin{tabular}{lc}
\toprule
\textbf{Registro} & \textbf{F1} \\
\midrule
207 & 0.1546 \\
108 & 0.7117 \\
113 & 0.7753 \\
\bottomrule
\end{tabular}
\caption{Registros mais cr\'iticos na detec\c{c}\~ao de picos R.}
\end{table}

\begin{figure}[H]
\centering
\includegraphics[width=\textwidth]{../figures/rpeak_detection_metrics.png}
\caption{Desempenho da detec\c{c}\~ao de picos R por registro.}
\end{figure}

Esse resultado sugere que o componente de detec\c{c}\~ao, em termos m\'edios, est\'a suficientemente est\'avel para sustentar an\'alises posteriores, mas h\'a registros espec\'ificos com comportamento patol\'ogico que precisam ser estudados individualmente.

\section{Interpreta\c{c}\~ao metodol\'ogica do estado atual}
At\'e o momento, o projeto j\'a estabeleceu uma base experimental coerente:

\begin{itemize}
    \item os dados foram reorganizados em pipeline reproduz\'ivel;
    \item a extra\c{c}\~ao de features por batimento j\'a est\'a operacional;
    \item o baseline cl\'assico j\'a fornece uma linha de compara\c{c}\~ao;
    \item a detec\c{c}\~ao de picos passou a ser avaliada explicitamente, e n\~ao assumida como perfeita.
\end{itemize}

Esse ponto \'e importante para o artigo futuro: a qualidade da etapa de segmenta\c{c}\~ao n\~ao pode ser tratada como detalhe de implementa\c{c}\~ao, mas como parte da validade interna do estudo.

\section{Limita\c{c}\~oes atuais}
Apesar de a base t\'ecnica estar bem encaminhada, h\'a limita\c{c}\~oes claras no estado atual:

\begin{itemize}
    \item o conjunto de atributos morfol\'ogicos ainda \'e simples;
    \item a etapa de split por registro ainda n\~ao explora valida\c{c}\~ao cruzada;
    \item os resultados de LSTM e CNN/LSTM ainda n\~ao est\~ao consolidados;
    \item a an\'alise dos registros problem\'aticos na detec\c{c}\~ao de picos ainda n\~ao foi formalizada;
    \item ainda n\~ao h\'a se\c{c}\~ao de revis\~ao bibliogr\'afica nem compara\c{c}\~ao com trabalhos relacionados.
\end{itemize}

\section{Como este texto pode evoluir para artigo}
O documento atual j\'a pode ser expandido para um manuscrito completo com a seguinte correspond\^encia:

\begin{longtable}{p{0.32\textwidth} p{0.60\textwidth}}
\toprule
\textbf{Se\c{c}\~ao futura do artigo} & \textbf{Base j\'a existente neste documento} \\
\midrule
Introdu\c{c}\~ao & problema de pesquisa, motiva\c{c}\~ao e objetivo geral; \\
Materiais e m\'etodos & base de dados, pipeline, features, modelo e m\'etricas; \\
Resultados & tabelas atuais de desempenho e figura da detec\c{c}\~ao de picos; \\
Discuss\~ao & interpreta\c{c}\~ao metodol\'ogica e limita\c{c}\~oes; \\
Conclus\~ao & consolida\c{c}\~ao dos achados e pr\'oximos passos. \\
\bottomrule
\end{longtable}

\section{Pr\'oximos passos recomendados}
Para aproximar o projeto de um artigo cient\'ifico robusto, os pr\'oximos passos mais importantes s\~ao:

\begin{enumerate}[label=\arabic*.]
    \item investigar qualitativamente os registros com pior detec\c{c}\~ao de picos, em especial o registro 207;
    \item expandir o conjunto de atributos morfol\'ogicos e temporais;
    \item consolidar experimentos com modelos recorrentes;
    \item acrescentar se\c{c}\~ao de trabalhos relacionados e justificar as escolhas metodol\'ogicas frente \`a literatura;
    \item registrar protocolos experimentais compar\'aveis entre abordagens cl\'assicas e profundas.
\end{enumerate}

\section{Comandos de reprodu\c{c}\~ao}
\subsection{Ambiente}
\begin{verbatim}
python -m venv .venv
source .venv/bin/activate
pip install --upgrade pip
pip install -r requirements.txt
pip install -r requirements-dev.txt
\end{verbatim}

\subsection{Pipeline principal}
\begin{verbatim}
python scripts/make_dataset.py
python scripts/preprocess_mitbih.py
python scripts/extract_symbolic_features.py
python scripts/train_classical_ml.py
\end{verbatim}

\subsection{Avalia\c{c}\~ao da detec\c{c}\~ao de picos}
\begin{verbatim}
MPLCONFIGDIR=/tmp/matplotlib .venv/bin/python scripts/evaluate_rpeak_detection.py
\end{verbatim}

\subsection{Testes}
\begin{verbatim}
pytest -q
\end{verbatim}

\subsection{Compila\c{c}\~ao deste documento}
\begin{verbatim}
cd reports/manuscript
pdflatex project_report.tex
pdflatex project_report.tex
\end{verbatim}

\end{document}
